% ============================================================
% 04_diseno_simulacion.tex
% Responsable: Integrante X
% ============================================================

\section{Diseño de la Simulación}

\subsection{Objetivo}

Implementar una simulación funcional de C-Lock en Python que permita:
\begin{itemize}
    \item Modelar múltiples cores accediendo a locks compartidos.
    \item Simular el mecanismo de clock-gating (suspensión/reanudación de cores).
    \item Medir el consumo energético simulado vs. un spin-lock clásico.
    \item Reproducir experimentos similares a los del paper.
\end{itemize}

\subsection{Componentes del simulador}

\begin{itemize}
    \item \textbf{\texttt{Core}}: Representa un core del procesador. Puede estar en estado
          \textit{running} o \textit{gated} (clock detenido).
    \item \textbf{\texttt{CLockManager}}: Gestiona el estado global de todos los locks y
          coordina los Conflict Checkers.
    \item \textbf{\texttt{Pool}}: Contiene los Conflict Checkers de un core.
    \item \textbf{\texttt{ConflictChecker}}: Detecta conflictos para un lock específico.
    \item \textbf{\texttt{Benchmark}}: Lanza múltiples cores ejecutando secciones críticas
          y mide métricas.
\end{itemize}

\subsection{Modelo de energía}

% TODO: Definir modelo de energía:
%   - Ciclos activos: costo = E_activo por ciclo
%   - Ciclos en clock-gating: costo = E_gated por ciclo (E_gated << E_activo)
%   - Comparar con spin-lock: core activo durante toda la espera

\subsection{Escenarios de prueba}

% TODO: Definir benchmarks a simular:
%   - N cores, 1 lock compartido (alta contención)
%   - N cores, M locks (baja contención)
%   - Comparar C-Lock vs spin-lock en ambos escenarios
